
\subsection{Quartic Scalar Coupling}\label{sec:CAGquartic}
A similar analysis can now be made for the quartic scalar coupling, with the general beta function 
\begin{equation}
    \beta_\lambda:=\mu \frac{\text{d} \alpha_\lambda}{\text{d}\mu}=\alpha_\lambda\left(d_1\alpha_\lambda + d_2\alpha_g+d_3\alpha_y\right)+d_4\alpha_g^2+d_5\alpha_y^2\;,\label{eq:betalam1loop}
\end{equation}
where $d_1,d_3,d_4\ge0$ and $d_2,d_5\le 0$. The quartic coupling, unlike the gauge beta function in Eq.\ \eqref{eq:betag1loopcaf} and the Yukawa beta function in Eq.\ \eqref{eq:betay1loop}, has terms that are not proportional to the coupling itself. Consequently, setting the quartic coupling to zero does not yield a trivial solution. 

Like the Yukawa and gauge couplings, we can first consider whether the quartic coupling can be asymptotically free on its own. In this case, the solution will be similar to that of the two other couplings
\begin{equation}
    \alpha_\lambda=\frac{\alpha_{\lambda_0}}{1-d_1 \alpha_{\lambda_0} \log \left(\frac{\mu}{\mu_0 }\right)}\;.
\end{equation}
The logarithmic term in the denominator is positive $d_1>0$, thus, the analysis is similar to that of the Yukawa coupling in isolation. At the energy scale $\mu_L=\mu_0\exp(1/{d_1\alpha_{\lambda_0}})$ the coupling will approach an UV Landau pole, while at energies $\mu<\mu_L$, the coupling is positive with a trivial Gaussian fixed point in the IR. Finally, above the Landau pole, $\mu>\mu_L$, the coupling is negative and runs to zero in the UV. We consider this a non-physical branch, since the potential would be unbounded from below. Consequently, the quartic coupling cannot exhibit asymptotic freedom without the gauge and Yukawa couplings. 

Let us now consider a theory with both gauge, Yukawa, and a quartic scalar coupling. The beta equation of the quartic coupling Eq.\ \eqref{eq:betalam1loop} can be solved using a similar approach to Eq.\ \eqref{eq:betayfrac} by taking the ratio $\frac{\beta_\lambda}{\beta_g}$ and we find
\begin{equation}
    \frac{\text{d} \alpha_\lambda}{\text{d}\alpha_g}=\frac{\alpha_\lambda}{b_0\alpha_g}\left(d_1\frac{\alpha_\lambda}{\alpha_g} + d_2+d_3\frac{\alpha_y}{\alpha_g}\right)+\frac{d_4}{b_0}+\frac{d_5 \alpha_y^2}{b_0\alpha_g^2}\;.
\end{equation}
To find the CAF conditions in a theory with a gauge, Yukawa, and quartic coupling, it is important to consider whether the Yukawa and gauge coupling are on or off fixed flow. This gives two different differential equations to solve and thus two sets of CAF conditions.

First, we consider the case of the Yukawa coupling off fixed flow. Then the differential equation simplifies to
\begin{equation}
    \frac{\text{d} \alpha_\lambda}{\text{d}\alpha_g}=\frac{\alpha_\lambda}{b_0\alpha_g}\left(d_1\frac{\alpha_\lambda}{\alpha_g} + d_2\right)+\frac{d_4}{b_0}\;.
    \label{eq:betalamfracapp}
\end{equation}
To find solutions, we will first define $k = (b_0 - d_2)^2 - 4d_1 d_4$, and $k_0 = (b_0 - d_2)\alpha_{g_0} - 2d_1\alpha_{\lambda0}$.
This will be useful for finding $\alpha_\lambda$ as a function of $\alpha_g$ for cases $k<0$, $k>0$, and $k=0$. 
\paragraph{Case 1 $k < 0$:}
We start with the general solution to the differential equation of $\alpha_\lambda(\alpha_g)$ (Equation \ref{eq:betalamfracapp}):
\begin{equation}
    \alpha_{\lambda} = \frac{\alpha_{g}}{2d_{1}} \left[ b_0 - d_2 + \sqrt{-k} \tan\left( \frac{\sqrt{-k}}{2b_0} \log\frac{\alpha_g}{\alpha_{g_0}} - \tan^{-1}\frac{k_0}{\sqrt{-k\alpha_{g_0}}} \right) \right] , \quad k<0 \,.
    \label{eq:alphalamkl0}
\end{equation}
This solution only holds for $k<0$. In this case the coupling will reach a Landau pole, because of the nature of the tangent function as one varies the gauge coupling. Thus, no possible CAF regions are found when $k<0$.

\paragraph{Case 2 $k > 0$:}
For $k>0$ the solution is given by
\begin{equation}
    \alpha_{\lambda} = \frac{\alpha_g}{2d_1} \left( b_0 - d_2 - \sqrt{k} \, \frac{ (k_0 + \sqrt{k}\alpha_{g_0}) + (k_0 - \sqrt{k}\alpha_{g_0}) \left( \frac{\alpha_g}{\alpha_{g_0}} \right)^{\sqrt{k}/b_0} }{ (k_0 + \sqrt{k}\alpha_{g_0}) - (k_0 - \sqrt{k}\alpha_{g_0}) \left( \frac{\alpha_g}{\alpha_{g_0}} \right)^{\sqrt{k}/b_0} } \right) , \quad k>0 \,,
    \label{eq:alphalamkg0}
\end{equation}
see Appendix \ref{app:caf} for derivation. 

Let us now do the CAF analysis for this case. From Eq.\ \eqref{eq:alphalamkg0}, it can be seen that there are two possible fixed flows between the gauge and self-coupling. Namely, when $k_0+\sqrt{k}\alpha_g=0$ or when $k_0-\sqrt{k}\alpha_g=0$ which gives us the expression
\begin{equation}
    \alpha_\lambda =\frac{b_0-d_2\pm\sqrt{k}}{2d_1}\alpha_g\;.\label{eq:lamgfixed}
\end{equation}
If we only consider solutions for $\alpha_\lambda>0$ to be relevant solutions, we will see that these fixed-flow lines also create the boundaries of the asymptotically free region of the quartic coupling. First, if $\frac{b_0-d_2-\sqrt{k}}{2d_1}>0$, then it is asymptotically free along both fixed-flow lines. If $\frac{b_0-d_2+\sqrt{k}}{2d_1}>0$ and $\frac{b_0-d_2-\sqrt{k}}{2d_1}<0$, it will only be asymptotically free along one direction. Lastly, if $\frac{b_0-d_2+\sqrt{k}}{2d_1}<0$, it will not be asymptotically free along any of the fixed-flow directions. 

To ensure that no Landau poles are encountered within the fixed-flow lines, we consider the denominator of Eq.\ \eqref{eq:alphalamkg0} to find other possible asymptotically free flows. Take the denominator to be zero, then the gauge coupling has to uphold
\begin{equation}
    \left(\frac{\alpha_g}{\alpha_{g_0}}\right)^{-\frac{\sqrt{k}}{b_0}}=\frac{k_0+\sqrt{k}\alpha_{g_0}}{k_0-\sqrt{k}\alpha_{g_0}}\;.
\end{equation}
Thus, if there are any solutions, there will be Landau poles. First, the solution of $k_0+\sqrt{k}\alpha_{g_0}=0$ is not relevant, because this zero will be cancelled by a zero in the numerator, as shown by the fixed-flow analysis. Thus, for $\alpha_g=0$ there are no Landau poles. However, for $\alpha_g>0$ there still might be poles. Since the LHS will be positive, so must the RHS, this means that both the numerator and denominator must have the same sign. Furthermore, we have $\alpha_g\to0$, ensuring that the LHS is less than one. Essentially, for a pole in the UV, the RHS will have to be less than one, and if the RHS is larger than one, we have a pole in the IR. Thus, for a pole in the UV
\begin{equation}
    0<\frac{k_0+\sqrt{k}\alpha_{g_0}}{k_0-\sqrt{k}\alpha_{g_0}}<1\;,
\end{equation}
For this to be the case, both the denominator and the numerator would have to be negative, for the numerator to be less than the denominator. Thus, to ensure no poles in the UV, we require that $k_0+\sqrt{k}\alpha_{g_0}>0$. Considering possible poles in the IR, the condition for a pole will be
\begin{equation}
    1<\frac{k_0+\sqrt{k}\alpha_{g_0}}{k_0-\sqrt{k}\alpha_{g_0}}\;.
\end{equation}
For this to be the case, we need both numerator and denominator to be positive. Since the numerator needs to be positive for no poles in the UV, we can ensure no poles in the IR by constraining the denominator $k_0-\sqrt{k}\alpha_{g_0}<0$. The reason why we need to have no poles in the IR, even though we are only interested in asymptotic freedom, is that the coupling would blow up at minus infinity, creating a potential which is unbounded from below. To ensure a bounded potential when having only one quartic coupling, we need the coupling to be positive at all RG scales. 

Thus, combining this analysis with Eq.\ \eqref{eq:lamgfixed} the restriction on $\alpha_{\lambda_0}$ and $\alpha_{g_0}$ is given by
\begin{equation}
    \frac{b_0-d_2-\sqrt{k}}{2d_1}\le \frac{\alpha_{\lambda_0}}{\alpha_{g_0}}\le \frac{b_0-d_2+\sqrt{k}}{2d_1}\;.
\end{equation}
We also need to ensure a positive quartic coupling at all scales, resulting in the condition of a positive upper boundary
\begin{equation}
    {b_0-d_2+\sqrt{k}}>0\;.
\end{equation}
Thus, the quartic coupling will be asymptotically free within the boundary of the fixed-flow lines and have at least one positive asymptotically free solution if
\begin{equation}
    k>0\;,\quad\quad  \frac{b_0-d_2-\sqrt{k}}{2d_1}\le \frac{\alpha_{\lambda_0}}{\alpha_{g_0}}\le \frac{b_0-d_2+\sqrt{k}}{2d_1}, \quad\quad {b_0-d_2+\sqrt{k}}>0\;.
\end{equation}


\paragraph{Case 3 $k = 0$:}
The solution is given by
\begin{equation}
    \alpha_{\lambda} = \frac{\alpha_g}{2d_1} \frac{ 4b_0 d_1 \alpha_{\lambda0} + k_0 (b_0 - d_2) \ln\frac{\alpha_g}{\alpha_{g_0}} }{ 2b_0 \alpha_{g_0} + k_0 \ln\frac{\alpha_g}{\alpha_{g_0}} } , \quad k=0 \,,
    \label{eq:alphalamkeq0}
\end{equation}
see Appendix \ref{app:caf} for derivation.

We have now obtained the last of the analytic solutions to Eq.\ \eqref{eq:betalamfracapp}, and we can move forward with the final CAF analysis. 

First, consider the denominator to be zero, this requirement yields
\begin{equation}
    \ln\frac{\alpha_g}{\alpha_{g_0}}=-\frac{2b_0\alpha_{g_0}}{k_0}\;,
\end{equation}
for $\alpha_g\ge0$ we again have two possible solutions, which will give Landau poles. For $\alpha_g<\alpha_{g_0}$, the RHS has to be negative, for a pole in the UV. To ensure no poles $k_0>0$. The situation is reversed for the IR, thus we require $k_0<0$. Inserting the only possible solution $(k_0=0)$ back into Eq.\ \eqref{eq:alphalamkeq0}, we obtain the restrictions for an asymptotically free quartic coupling in this case to be
\begin{equation}
    \alpha_{\lambda} = \frac{ \alpha_{\lambda0}  }{ \alpha_{g_0}  } \alpha_g\;,\quad\quad k_0=0\;, \quad\quad k=0\;.
\end{equation}
Thus for $k=0$, only one flow line would be asymptotically free. 

We have now considered the case of the gauge and Yukawa coupling off fixed flow. However, for the gauge and Yukawa coupling on fixed flow, the ratio of the beta function of the quartic and gauge coupling does not reduce
\begin{equation}
    \frac{\text{d} \alpha_\lambda}{\text{d}\alpha_g}=\frac{\alpha_\lambda}{b_0\alpha_g}\left(d_1\frac{\alpha_\lambda}{\alpha_g} + d_2+d_3\frac{\alpha_y}{\alpha_g}\right)+\frac{d_4}{b_0}+\frac{d_5\alpha_y^2}{b_0\alpha_g^2}\;.
\end{equation}
The CAF conditions in this case can be found by removing the Yukawa dependency using the fixed-flow relation of the Yukawa and gauge coupling in Eq.\ \eqref{eq:ygfixed2}. This lets us replace $\frac{\alpha_y}{\alpha_g}=\frac{b_0-c_1}{c_2}$ 
\begin{equation}
    \frac{\text{d} \alpha_\lambda}{\text{d}\alpha_g}=\frac{\alpha_\lambda}{b_0\alpha_g}\left(d_1\frac{\alpha_\lambda}{\alpha_g} + d_2+d_3\frac{b_0-c_1}{c_2}\right)+\frac{d_4}{b_0}+\frac{d_5(b_0-c_1)^2}{b_0c_2^2}\;.
\end{equation}
Now the equation can be rewritten as
\begin{equation}
    \frac{\text{d} \alpha_\lambda}{\text{d}\alpha_g}=\frac{\alpha_\lambda}{b_0\alpha_g}\left(d_1\frac{\alpha_\lambda}{\alpha_g} + d_2'\right)+\frac{d_4'}{b_0}\;,
\end{equation}
where 
\begin{align}
    d_2'=d_2+d_3\frac{b_0-c_1}{c_2}\;,\quad\quad d_4'&={d_4}+d_5\frac{(b_0-c_1)^2}{c_2^2}\;.
\end{align}

Now that we have done the full CAF analysis, we can briefly sum up the results: The conditions when the Yukawa and gauge couplings are not on fixed flow (derived in Eq.\ \eqref{eq:CAF1con}) are found to be
\begin{equation}
    b_0<0,\quad\quad b_0-c_1>0,\quad\quad k\geq 0,\quad\quad b_0-d_2+\sqrt{k}>0\;,\label{eq:CAF}
\end{equation}
where $k=(b_0-d_2)^2-4d_1d_4$. Additionally, the constraints on the gauge, Yukawa, and quartic coupling at the reference energy scale (derived in Eqs. \eqref{eq:ke0con} and \eqref{eq:kg0con}) are given by
\begin{align}
    &\text{For }k=0: \quad \alpha_{\lambda} = \frac{ \alpha_{\lambda0}  }{ \alpha_{g_0}  } \alpha_g\;,\quad\quad k_0=0\;,\label{eq:CAFalpha}\\&
    \text{For }k>0: \quad 
    \frac{b_0-d_2-\sqrt{k}}{2d_1}\le \frac{\alpha_{\lambda_0}}{\alpha_{g_0}}\le \frac{b_0-d_2+\sqrt{k}}{2d_1}\;, \label{eq:CAFalphafixed}
\end{align}
where $k_0 = (b_0 - d_2)\alpha_{g_0} - 2d_1\alpha_{\lambda0}$. 

With the CAF conditions off fixed flow out of the way, the CAF conditions on fixed flow (see Eq.\ \eqref{eq:appCAF1conon}) are given by
\begin{align}
    &b_0<0,\quad\quad b_0-c_1>0,\quad\quad k'\geq 0,\quad \quad  b_0-d_2'+\sqrt{k'}>0\;,\label{eq:CAFfixed}
\end{align}  
where 
\begin{align}
    d_2'&=d_2+d_3\frac{b_0-c_1}{c_2}\;,\quad
    d_4'=d_4+d_5\left(\frac{b_0-c_1}{c_2}\right)^2\;,\quad
    k'= \left(b_0-d_2'\right)^2-4d_1d_4'\;.
\end{align}
It is important to note that for the CAF conditions on fixed flow, the primed coefficients can change sign. This point was missing from the analysis in \cite{Pica:2016krb}. As noted above, the $d$ coefficients satisfy $d_2\le0$ and $d_4\ge0$. However, on fixed flow, this will change since we obtain
\begin{align}
     &d_2'>0 \quad \text{  if}\quad d_3\frac{b_0-c_1}{c_2}>|d_2|\;,
     \\&d_4'<0 \quad \text{  if}\quad |d_5|\frac{(b_0-c_1)^2}{c_2^2}>d_4\;.
\end{align}
A sign change on the primed coefficients can, in some cases, ensure that $k'$ and $b_0-d_2'+\sqrt k$ are positive for all parameter values. Consequently, only the first two conditions may restrict the possibility of the theory exhibiting CAF. We will see the effect of this in later sections.

Similar to the case of the Yukawa coupling, we can visualise the quartic and gauge coupling off- and on-fixed-flow regimes. This can be seen in Figure~\ref{fig:glamCAFplot}, where values are shown, such that the condition in Eq. \eqref{label} are upheld. Note that this plot is for the case of the gauge and Yukawa coupling off fixed flow. Again the fixed-flow condition, this time between the quartic and the gauge coupling, is more restrictive in the values of the couplings at some reference scale, since they only occur along the two green trajectories. In between these the couplings will be asymptotically free, and a CAF regime can be realised.  

The analysis has thus far been restricted to simple scenarios involving either a single gauge, Yukawa, or quartic scalar coupling or combinations thereof. We now generalise this approach to the case of multiple couplings, where, as it will turn out, the CAF conditions cannot always be found analytically from the coupled differential equations. 